% F012 双根号函数值域：几何意义法（距离型）
\begin{statementbox}
\textbf{\textcolor{CStatement}{它解决什么问题？}}

遇到两个根号相加，例如
\[
  \sqrt{(x-1)^2+4}+\sqrt{(x-3)^2+9},
\]
想知道它最小能等于多少，或者所有可能的值有哪些。一个一个拆根号很难，直接求导也不够省事。本结论把两个根号变成\textbf{两段路程}，利用\textbf{两点之间线段最短}解决。

\medskip
\textbf{\textcolor{CStatement}{使用条件}}
\[
  f(x)=\sqrt{(x-m)^2+a^2}+\sqrt{(x-n)^2+b^2},
\]
其中 $a>0$，$b>0$，$m,n$ 是固定的实数，$x$ 可以取任意实数。
\begin{center}\begin{tikzpicture}[font=\small,>=stealth,line join=round]
\draw[gray!40,rounded corners=3pt] (0,0) rectangle (14.2,8.6);
\node at (7.1,8.0) {两个根号是两段路：折线拉直，路程最短};
\begin{scope}[shift={(2.5529411764705885,4.25)},x=0.8235294117647058cm,y=0.8235294117647058cm]
\draw[->,gray] (-2,0)--(4,0) node[right]{$x$};
\draw[blue,very thick] (1,-2)--(3,3);
\draw[orange!85!black,dashed,thick] (1,-2)--(-1,0)--(3,3);
\draw[dotted,gray] (1,0)--(1,-2) (3,0)--(3,3);
\fill (1,-2) circle(2pt) (3,3) circle(2pt);
\fill[blue] (1.8,0) circle(2.3pt);\fill[orange!85!black] (-1,0) circle(2pt);
\node[below right=5pt] at (1,-2) {$A$};\node[above right=5pt] at (3,3) {$B$};
\node[below right=8pt] at (1.8,0) {$P_0$};\node[above left=6pt] at (-1,0) {$P$};
\end{scope}
\node[align=left,anchor=west] at (8.2,4.9) {$A=(1,-2)$\\[9pt]$B=(3,3)$\\[12pt]$P_0=(9/5,0)$\\[12pt]$AP_0+P_0B=AB$\\[12pt]最短路程 $\sqrt{29}$};
\node[text width=12.8cm,align=center] at (7.1,.65) {横纵坐标等比例；蓝线为最短直路，橙色虚线为经其他位置 $P$ 的路。};
\end{tikzpicture}\end{center}



\medskip
\end{statementbox}
\clearpage
\begin{statementbox}
\textbf{\textcolor{CStatement}{核心结论}}

取三个点
\[
 A(m,-a),\qquad B(n,b),\qquad P(x,0).
\]
由两点间距离公式，
\[
  f(x)=|PA|+|PB|\geq |AB|
  =\boxed{\sqrt{(m-n)^2+(a+b)^2}}.
\]
因为 $A$ 在 $x$ 轴下方、$B$ 在上方，线段 $AB$ 一定穿过 $x$ 轴。让 $P$ 正好落在交点上，两段路就连成了一段直线，所以等号\textbf{一定能够成立}。

最小值恰在
\[
 \boxed{x_0=\frac{bm+an}{a+b}}
\]
处取得。因此完整值域为
\[
 \boxed{\left[\sqrt{(m-n)^2+(a+b)^2},\ +\infty\right)}.
\]

\textbf{一句话记忆：} 两个根号是两段距离；把定点放在轴的两侧，折线拉直就是最短。

\medskip
\textbf{\textcolor{CStatement}{注意：}} 上面的完整值域以 $x\in\mathbb R$ 为前提。若题目限制了 $x$，必须重新检查最短位置 $x_0$ 是否允许取到。
\end{statementbox}
